\textit{Magnificat en re major} (BWV 243) és una de les grans obres vocals de Johann
Sebastian Bach, publicada el 1723 i escrita per a cor a cinc veus i orquestra. La corda
d'un violí fa 32,5 cm de llargària i la freqüència fonamental corresponent al re major és
de 38,89 Hz.

\begin{apartat}{1,25}
Dibuixeu l'harmònic fonamental i el tercer harmònic i indiqueu-ne els nodes i els ventres.
Determineu les longituds d'ona de cadascun dels harmònics i calculeu la velocitat de
propagació de l'ona que produeix aquesta nota en la corda del violí. Quina serà la
freqüència de l'harmònic fonamental si reduïm la llargària de la corda a 27 cm tot
mantenint la mateixa tensió a la corda?
\end{apartat}

\begin{apartat}{1,25}
Un espectador situat al segon pis del Palau de la Música Catalana percep un nivell
d'intensitat sonora de 30 dB, corresponent al so dels violins, situats a 23 m. Quina és la
potència amb què els violins emeten el so? Si, degut a les restriccions per la COVID-19,
el nombre de violinistes es redueix a la meitat, quin serà el nivell de la intensitat
sonora generada pels violins a 23 m de distància?
\end{apartat}

\dades{$I_0 = 10^{-12}\ \mathrm{W\,m^{-2}}$.}
